\begin{abstract}
Conventional Hyperspectral Imaging (HSI) is based on acquiring either reflectance, emission, or excitation profiles from each pixel of an image. 
Emission-based HSI can be expanded into a fourth dimension by varying the excitation wavelength, which enables an improved discrimination of materials that contain mixed fluorophores. 
Multi-excitation HSI (ME-HSI) creates a 4D dataset encompassing signal intensity across two spatial and two spectral dimensions, the latter being the emission and excitation wavelengths. 
The processing of ME-HSI datasets present notable challenges due to the numerous spectral layers per pixel, introducing considerable redundancy and computational burden. 
Existing band selection methods, developed for 3D HSI, fail to capture the nonlinear cross-excitation correlations present in ME-HSI. 
To fill this gap, we developed a deep learning framework comprising three stages: (1) a 3D convolutional autoencoder with parallel excitation branches that learns a unified latent representation of ME-HSI data, (2) a perturbation-based attribution method that traces reconstruction sensitivity to individual excitation-emission wavelength pairs, and (3) Maximum Marginal Relevance (MMR) selection that balances informativeness with spectral diversity.
Evaluated on lichen samples with four ground-truth classes across 3,072 configurations, the framework achieved the accuracy exceeding the baseline, with 95\% data reduction, which corresponds to only 9 out of 192 bands.
The presented framework facilitates practical applications of ME-HSI by reducing acquisition complexity while improving discriminative power.
\end{abstract}